%% kapa: exam=1 round=32 year=2021 period=1 form=A subject=real_estate subject_ko=부동산학원론 seq=24 number=104
%% source: Q-Net 감정평가사 기출문제 (cst00309) artlSeq=5208397 file=extracted/1008/5208397/A형/2021년도 제32회 감정평가사 1차 1교시 A형.pdf
%% conversion: kordoc 4.24.51 -> tools/gen1.py (auto); review=proofread-sonnet; flags=none
\begin{question}{104}
조기상환에 관한 설명으로 옳지 않은 것은?
\begin{choices}
  \item 조기상환이 어느 정도 일어나는가를 측정하는 지표로 조기상환율(CPR: Constant Prepayment Rate)이 있다.
  \item 저당대출차입자에게 주어진 조기상환권은 풋옵션(put option)의 일종으로, 차입자가 조기상환을 한다는 것은 대출잔액을 행사가격으로 하여 대출채권을 매각하는 것과 같다.
  \item 저당대출차입자의 조기상환 정도에 따라 MPTS(Mortgage Pass-Through Securities)의 현금흐름과 가치가 달라진다.
  \item 이자율 하락에 따른 위험을 감안하여 금융기관은 대출기간 중 조기상환을 금지하는 기간을 설정하고, 위반 시에는 위약금으로 조기상환수수료를 부과하기도 한다.
  \item 저당대출차입자의 조기상환은 MPTS(Mortgage Pass-Through Securities) 투자자에게 재투자 위험을 유발한다.
\end{choices}
\end{question}
\answer{2}
